
El ciclo de vida del proyecto define cómo se organiza el trabajo desde su inicio hasta el cierre, 
estableciendo las fases que producen entregables verificables y que sirven como base para 
planificar el alcance, controlar cambios y validar resultados. La metodología de desarrollo de 
software adoptada debe respetarlo, adaptándose internamente según la naturaleza del producto.

\subsection{Ciclo de vida del proyecto}

El proyecto seguirá el ciclo de vida estándar estructurado en cinco fases fundamentales, tal 
como lo establece el marco de gestión de proyectos de referencia. Estas fases son universales 
al tipo de proyecto y aplican independientemente de la metodología técnica de desarrollo de 
software que se adopte internamente:

\begin{itemize}
    \item \textbf{Inicio:} Se formaliza la existencia del proyecto, se define el problema central 
    (ineficiencia en el trámite de licencias de conducción a nivel nacional), se identifican los 
    stakeholders principales y se establece la autorización para planificar. Esta fase produce el 
    acta de constitución del proyecto y el registro inicial de interesados.

    \item \textbf{Planificación:} Se elabora el Plan de Gestión del Alcance, incluyendo los entregables, exclusiones, restricciones, supuestos, la 
    EDT/WBS y la línea base del proyecto. En esta fase se define qué se construirá, quién lo 
    construirá y bajo qué condiciones. Esta fase puede refinarse progresivamente conforme 
    avanza el conocimiento del proyecto.

    \item \textbf{Ejecución:} Se construye el sistema de software mediante el desarrollo 
    iterativo e incremental de sus módulos funcionales. El equipo implementa los diseños 
    aprobados, gestiona integraciones externas y produce los 
    incrementos verificables del producto.

    \item \textbf{Monitoreo y Control:} A lo largo de todo el ciclo de vida se realiza seguimiento 
    al avance, se controlan cambios al alcance mediante el proceso formal de control integrado 
    de cambios y se validan los entregables producidos en cada iteración. Esta fase es 
    transversal a todas las demás.

    \item \textbf{Cierre:} Se formalizan las actas de aceptación de los entregables finales, se 
    consolidan las lecciones aprendidas, se cierra el contrato con proveedores y se entrega 
    oficialmente el sistema al cliente.
\end{itemize}

\subsection{Metodología de desarrollo de software}

Para el desarrollo del sistema de trámites de licencias, se adopta una metodología ágil incremental basada en el marco de trabajo Scrum. 

La elección de esta metodología se fundamenta en que el sistema está compuesto por múltiples módulos funcionales independientes. En lugar de utilizar un enfoque tradicional (cascada) que entrega todo al final, Scrum permite construir el software mediante entregas parciales (incrementos), aportando valor continuo a la entidad de tránsito y permitiendo la adaptación ante posibles cambios en integraciones complejas (como RUNT o pasarelas de pago).

Para que esta metodología deje de ser un concepto general y se aplique operativamente al desarrollo, se define la siguiente estructura de trabajo:

\begin{itemize}
    \item \textbf{Ciclos de trabajo (Sprints):} La fase de ejecución se divide en iteraciones de tiempo fijo (Sprints de 2 a 4 semanas). En cada Sprint, el equipo aborda un conjunto específico de requisitos de la EDT (por ejemplo, el módulo de turnos) y ejecuta su análisis, diseño, codificación y pruebas, dando como resultado un incremento de software 100\% funcional.
    
    \item \textbf{Roles de la metodología:} La ejecución técnica estará liderada por perfiles ágiles integrados en la estructura organizacional:
    \begin{itemize}
        \item \textit{Product Owner:} Representado por el Analista Funcional en coordinación con el Director de Proyecto del Cliente (Ministerio de Transporte). Son los encargados de priorizar qué se desarrolla primero.
        \item \textit{Scrum Master / Project Manager:} Encargado de liderar la metodología, eliminar bloqueos e integrar áreas transversales (Infraestructura, Seguridad).
        \item \textit{Equipo de Desarrollo:} Compuesto por Arquitecto de Software, desarrolladores (Front/Back), QA y DevOps.
    \end{itemize}
    
    \item \textbf{Ceremonias (Eventos de validación):} El trabajo incremental se controla mediante reuniones formales. Diariamente el equipo operativo realiza una Daily Scrum para sincronizar tareas. Al final de cada Sprint, se ejecuta una Sprint Review junto a los directivos y supervisores del cliente para validar y aceptar el módulo entregado.
\end{itemize}

\subsection{Adaptaciones del ciclo de vida para el proyecto de software}

Dado que el marco general de gestión del proyecto exige control y validación de entregables, Scrum se adapta al ciclo de vida del proyecto bajo los siguientes principios:

\subsubsection{Adaptación A-01: División de la fase de ejecución en Sprints}
La fase de ejecución tradicional se adapta para no ser un bloque secuencial único. En su lugar, el sistema se desarrolla iterativamente. 
\begin{itemize}
    \item \textit{Impacto en la gestión:} Desarrollar por Sprints mitiga el riesgo técnico de las fallas de integración. Cada Sprint tiene su propia acta de entrega parcial y debe cumplir con los criterios de aceptación documentados (ver Sección \ref{sec:Entregables del producto}).
    \item \textit{Responsable:} Project Manager / Scrum Master.
\end{itemize}

\subsubsection{Adaptación A-02: Planificación progresiva (Gestión del Backlog)}
La EDT/WBS se trata como un documento vivo (Product Backlog). Al inicio de cada Sprint, la planificación del alcance se refina progresivamente.
\begin{itemize}
    \item \textit{Impacto en la gestión:} Debido a la naturaleza normativa de las licencias virtuales y biometría, es normal que los detalles técnicos se aclaren durante el desarrollo. Sin embargo, cualquier modificación que afecte la línea base del proyecto debe pasar por un proceso formal de control de cambios.
    \item \textit{Responsable:} Product Owner / Analista Funcional.
\end{itemize}

\subsubsection{Adaptación A-03: Validación del alcance por Incremento}
En lugar de una única validación al final del proyecto, cada incremento funcional producido en un Sprint es sometido a un proceso de validación formal en la Sprint Review.
\begin{itemize}
    \item \textit{Impacto en la gestión:} Esta demostración funcional permite al representante del cliente aprobar los módulos entregados. Las actas de aceptación parcial de cada Sprint sirven como insumos obligatorios para el acta de cierre definitivo del proyecto.
    \item \textit{Responsable:} Project Manager junto al Director de Proyecto del Cliente.
\end{itemize}

